\section{Programowanie matematyczne i algorytmy zaawansowane}

\subsection{Metody poszukiwania ekstremum funkcji nieliniowej}

Rozważamy problem
\[
\min_{x\in\R^n} f(x)
\]
albo równoważnie maksymalizację przez minimalizację $-f(x)$. Funkcja może być wypukła albo niewypukła, gładka albo niegładka, tania albo kosztowna w ewaluacji. Wybór metody zależy głównie od dostępności pochodnych, wymiaru, wypukłości, szumu i kosztu obliczenia funkcji.

\subsubsection{Warunki optymalności bez ograniczeń}

Jeśli $f$ jest różniczkowalna i $x^*$ jest lokalnym minimum wewnętrznym, to konieczny warunek pierwszego rzędu:
\[
\nabla f(x^*)=0.
\]
Jeśli $f$ jest dwa razy różniczkowalna, konieczny warunek drugiego rzędu:
\[
\nabla^2 f(x^*) \succeq 0.
\]
Warunek wystarczający dla ścisłego minimum lokalnego:
\[
\nabla f(x^*)=0,\qquad \nabla^2 f(x^*)\succ 0.
\]
Dla funkcji wypukłej każdy punkt spełniający $\nabla f(x^*)=0$ jest minimum globalnym. Dla funkcji ściśle wypukłej minimum globalne jest co najwyżej jedno.

\subsubsection{Metoda największego spadku}

Metoda największego spadku, czyli gradient descent, wykonuje kroki:
\[
x_{k+1}=x_k-\alpha_k\nabla f(x_k),
\]
gdzie $\alpha_k>0$ jest długością kroku. Kierunek $-\nabla f(x_k)$ jest kierunkiem najszybszego lokalnego spadku w normie euklidesowej.

Długość kroku można dobrać:
\begin{itemize}
  \item jako stałą,
  \item przez dokładne przeszukiwanie liniowe,
  \item przez warunki Armijo albo Wolfe'a,
  \item według harmonogramu malejącego, zwłaszcza w metodach stochastycznych.
\end{itemize}

Dla funkcji wypukłej z gradientem Lipschitza, przy odpowiednim kroku metoda zbiega do minimum. Dla funkcji silnie wypukłej zbieżność jest liniowa, ale może być bardzo wolna przy złym uwarunkowaniu, gdy poziomice są wydłużone.

\subsubsection{Metoda Newtona}

Metoda Newtona używa przybliżenia kwadratowego:
\[
f(x_k+p)\approx f(x_k)+\nabla f(x_k)^Tp+\frac{1}{2}p^T\nabla^2f(x_k)p.
\]
Kierunek Newtona spełnia:
\[
\nabla^2 f(x_k)p_k=-\nabla f(x_k),
\]
a aktualizacja:
\[
x_{k+1}=x_k+\alpha_k p_k.
\]

Jeśli Hessian jest dodatnio określony i punkt startowy jest dostatecznie blisko minimum, metoda ma zbieżność kwadratową. Wady: koszt obliczenia i odwracania/rozwiązywania układu z Hessianem, problemy przy Hessianie osobliwym albo nieokreślonym oraz brak globalnych gwarancji bez line search albo trust region.

\subsubsection{Metody quasi-Newtonowskie}

Metody quasi-Newtonowskie przybliżają Hessian albo jego odwrotność na podstawie kolejnych gradientów. Najbardziej znana jest BFGS. Niech
\[
s_k=x_{k+1}-x_k,\qquad y_k=\nabla f(x_{k+1})-\nabla f(x_k).
\]
BFGS aktualizuje przybliżenie Hessianu tak, aby spełniało równanie secant:
\[
B_{k+1}s_k=y_k.
\]
L-BFGS przechowuje tylko kilka ostatnich par $(s_k,y_k)$, dlatego nadaje się do dużych wymiarów.

Zalety: zwykle szybsze niż gradient descent, bez jawnego Hessianu. Wady: wymagają gradientu, a dla funkcji silnie niewypukłych lub zaszumionych mogą wymagać zabezpieczeń.

\subsubsection{Metody bezpochodne}

Gdy gradient nie jest dostępny, funkcja jest szumowa albo niedokładna, stosuje się metody bezpochodne:
\begin{description}
  \item[Nelder-Mead] utrzymuje sympleks $n+1$ punktów i wykonuje odbicie, ekspansję, kontrakcję i redukcję. Działa dobrze w małych wymiarach, ale nie ma silnych gwarancji w ogólności.
  \item[Pattern search] bada skończony zbiór kierunków i zmniejsza krok, gdy nie ma poprawy.
  \item[CMA-ES] adaptuje rozkład normalny próbkowania i macierz kowariancji. Jest silny dla trudnych, niewypukłych problemów ciągłych, ale kosztowny.
  \item[Symulowane wyżarzanie] akceptuje czasem gorsze ruchy z prawdopodobieństwem malejącym wraz z temperaturą, co pomaga wyjść z minimów lokalnych.
\end{description}

\subsubsection{Optymalizacja stochastyczna}

Dla funkcji będącej średnią po danych:
\[
f(x)=\frac{1}{m}\sum_{i=1}^{m}\ell_i(x)
\]
pełny gradient może być kosztowny. Stochastic Gradient Descent używa losowej próbki albo mini-batcha:
\[
x_{k+1}=x_k-\alpha_k g_k,\qquad \mathbb{E}[g_k]=\nabla f(x_k).
\]
W uczeniu maszynowym popularne są modyfikacje z momentum, RMSProp, Adam i AdamW. Ich zaletą jest skalowalność, wadą wrażliwość na hiperparametry i trudniejsza analiza końcowej dokładności.

\subsubsection{Lokalne i globalne optimum}

Dla funkcji niewypukłych metoda lokalna zwykle gwarantuje co najwyżej znalezienie punktu stacjonarnego albo lokalnego minimum. Globalna optymalizacja wymaga dodatkowych technik: wielostartu, branch and bound, relaksacji wypukłych, metaheurystyk, metod przedziałowych albo specjalnej struktury problemu.

\subsubsection{Co powiedzieć na obronie}

Najważniejsze metody bez ograniczeń to gradient descent, Newton, quasi-Newton, metody bezpochodne i stochastyczne. Gradient descent jest prosty, ale może być wolny; Newton jest szybki lokalnie, ale kosztowny; BFGS/L-BFGS to praktyczny kompromis; Nelder-Mead i inne metody bezpochodne są użyteczne, gdy nie mamy gradientu. W problemach wypukłych optimum lokalne jest globalne, w niewypukłych nie.

\subsection{Metody poszukiwania ekstremum w obecności ograniczeń}

Problem z ograniczeniami ma postać:
\[
\min_x f(x)
\]
przy warunkach
\[
g_i(x)\le 0,\quad i=1,\dots,m,\qquad h_j(x)=0,\quad j=1,\dots,p.
\]
Zbiór dopuszczalny jest zbiorem punktów spełniających ograniczenia. Ograniczenia mogą być liniowe, nieliniowe, wypukłe, niewypukłe, równościowe, nierównościowe albo całkowitoliczbowe.

\subsubsection{Warunki KKT}

Dla problemu gładkiego definiujemy Lagrangian:
\[
\mathcal{L}(x,\lambda,\mu)=f(x)+\sum_{i=1}^{m}\lambda_i g_i(x)+\sum_{j=1}^{p}\mu_j h_j(x),
\]
gdzie $\lambda_i\ge 0$ są mnożnikami dla nierówności, a $\mu_j$ dla równości.

Warunki Karusha-Kuhna-Tuckera:
\begin{align*}
\nabla_x \mathcal{L}(x^*,\lambda^*,\mu^*)&=0 &&\text{stacjonarność},\\
g_i(x^*)&\le 0,\quad h_j(x^*)=0 &&\text{dopuszczalność pierwotna},\\
\lambda_i^*&\ge 0 &&\text{dopuszczalność dualna},\\
\lambda_i^* g_i(x^*)&=0 &&\text{komplementarność}.
\end{align*}
Dla problemów wypukłych przy odpowiednich warunkach regularności KKT są nie tylko konieczne, ale też wystarczające.

\subsubsection{Metoda mnożników Lagrange'a}

Dla samych równości $h_j(x)=0$ szukamy punktów stacjonarnych Lagrangianu:
\[
\nabla f(x^*)+\sum_{j=1}^{p}\mu_j\nabla h_j(x^*)=0,\qquad h_j(x^*)=0.
\]
Geometrycznie gradient funkcji celu w optimum leży w rozpięciu gradientów aktywnych ograniczeń, czyli nie ma składowej pozwalającej poprawić wynik wzdłuż zbioru dopuszczalnego.

\subsubsection{Metody kar i barier}

Metody kar zamieniają problem ograniczony na sekwencję problemów bez ograniczeń:
\[
\min_x f(x)+\rho_k P(x),
\]
gdzie $P(x)$ karze naruszenie ograniczeń, np.
\[
P(x)=\sum_i \max(0,g_i(x))^2+\sum_j h_j(x)^2.
\]
Gdy $\rho_k$ rośnie, rozwiązania powinny zbliżać się do dopuszczalnych. Wadą jest pogarszające się uwarunkowanie dla dużych kar.

Metody barier działają wewnątrz zbioru dopuszczalnego, np. dla $g_i(x)<0$:
\[
\min_x f(x)-\mu\sum_i \log(-g_i(x)).
\]
Gdy $\mu\to 0$, bariera słabnie, a rozwiązanie zbliża się do optimum problemu oryginalnego. To idea metod punktu wewnętrznego.

\subsubsection{Programowanie sekwencyjne kwadratowe}

Sequential Quadratic Programming rozwiązuje w każdej iteracji przybliżony problem kwadratowy: funkcję celu zastępuje modelem kwadratowym Lagrangianu, a ograniczenia linearyzuje. SQP jest bardzo skuteczne dla gładkich problemów nieliniowych średniego rozmiaru.

\subsubsection{Metody projekcyjne}

Dla prostego zbioru dopuszczalnego $C$ można wykonać krok gradientowy i projekcję:
\[
x_{k+1}=P_C(x_k-\alpha_k\nabla f(x_k)),
\]
gdzie
\[
P_C(y)=\argmin_{x\in C}\|x-y\|_2.
\]
Metoda jest naturalna dla ograniczeń pudełkowych, kul, sympleksu albo prostych zbiorów wypukłych. W uczeniu maszynowym odpowiada np. obcinaniu wag do zakresu albo projekcji na ograniczenie normy.

\subsubsection{Active set i interior point}

Metody aktywnego zbioru zgadują, które ograniczenia nierównościowe są aktywne w optimum, czyli spełnione jako równości. Następnie rozwiązują problem z tymi ograniczeniami jak z równościami i aktualizują zbiór aktywny.

Metody punktu wewnętrznego poruszają się we wnętrzu zbioru dopuszczalnego i używają barier. Są podstawą bardzo wydajnych solverów dla programowania liniowego, kwadratowego i wypukłego.

\subsubsection{Ograniczenia całkowitoliczbowe}

Jeśli część zmiennych musi być całkowita albo binarna, problem staje się programowaniem całkowitoliczbowym. Typowe techniki:
\begin{itemize}
  \item branch and bound,
  \item branch and cut,
  \item relaksacja liniowa albo wypukła,
  \item płaszczyzny tnące,
  \item heurystyki naprawcze.
\end{itemize}
Tego typu problemy są często NP-trudne.

\subsubsection{Co powiedzieć na obronie}

Dla ograniczeń kluczowe są warunki KKT: stacjonarność, dopuszczalność, nieujemność mnożników i komplementarność. Metody praktyczne to mnożniki Lagrange'a, kary, bariery, projekcja, SQP, active set i interior point. Dla problemów wypukłych warunki KKT dają globalną charakterystykę optimum; dla niewypukłych zwykle tylko lokalną.

\subsection{Wielomianowy schemat aproksymacyjny}

Wiele problemów NP-trudnych jest zbyt kosztownych do dokładnego rozwiązania dla dużych instancji. Algorytmy aproksymacyjne zwracają rozwiązania bliskie optymalnym z formalną gwarancją jakości.

\subsubsection{Współczynnik aproksymacji}

Dla problemu minimalizacyjnego algorytm ma współczynnik $\rho(n)\ge 1$, jeśli dla każdej instancji:
\[
A(I)\le \rho(n)\cdot OPT(I),
\]
gdzie $A(I)$ jest wartością rozwiązania algorytmu, a $OPT(I)$ optimum. Dla problemu maksymalizacyjnego zwykle wymaga się:
\[
A(I)\ge \frac{1}{\rho(n)}OPT(I).
\]
Jeśli $\rho$ jest stałe, mówimy o stałej aproksymacji.

\subsubsection{PTAS}

Polynomial-Time Approximation Scheme to rodzina algorytmów, która dla dowolnego $\eps>0$ zwraca rozwiązanie o jakości:
\[
A(I)\le (1+\eps)OPT(I)
\]
dla minimalizacji albo
\[
A(I)\ge (1-\eps)OPT(I)
\]
dla maksymalizacji, w czasie wielomianowym względem rozmiaru wejścia $n$ dla ustalonego $\eps$.

Istotne: czas może być bardzo zły jako funkcja $1/\eps$, np.
\[
\BigO(n^{1/\eps}).
\]
Dla każdego ustalonego $\eps$ jest to wielomian w $n$, więc formalnie PTAS.

\subsubsection{FPTAS}

Fully Polynomial-Time Approximation Scheme wymaga czasu wielomianowego zarówno względem $n$, jak i $1/\eps$, np.
\[
\BigO(n^3/\eps^2).
\]
FPTAS jest znacznie silniejszy praktycznie i teoretycznie niż PTAS.

\subsubsection{EPTAS}

Efficient PTAS ma czas postaci
\[
f(1/\eps)\cdot n^{O(1)},
\]
gdzie wykładnik przy $n$ nie zależy od $\eps$. To kompromis między PTAS i FPTAS.

\subsubsection{Przykłady}

\begin{itemize}
  \item Problem plecakowy 0/1 ma FPTAS oparty na skalowaniu wartości przedmiotów i programowaniu dynamicznym.
  \item Euklidesowy TSP w ustalonym wymiarze ma PTAS.
  \item Dla ogólnego metrycznego TSP znany jest algorytm Christofidesa z gwarancją $3/2$, ale nie jest to PTAS.
  \item Dla wielu problemów silnie NP-trudnych istnienie FPTAS implikowałoby $\mathrm{P}=\NP$.
\end{itemize}

\subsubsection{Schemat FPTAS dla plecaka}

W plecaku maksymalizujemy wartość przy ograniczeniu wagi. Klasyczne DP po wartościach działa w czasie zależnym od sumy wartości, która może być pseudowielomianowa. FPTAS skaluje wartości:
\[
v_i'=\left\lfloor \frac{v_i}{K}\right\rfloor
\]
dla odpowiednio dobranego $K$ zależnego od $\eps$ i $v_{\max}$. Następnie rozwiązujemy problem DP dla mniejszych wartości. Tracimy co najwyżej kontrolowany ułamek optimum.

\subsubsection{Co powiedzieć na obronie}

PTAS to rodzina algorytmów, która dla dowolnej dokładności $\eps$ daje rozwiązanie w granicach $1+\eps$ od optimum w czasie wielomianowym dla ustalonego $\eps$. FPTAS jest silniejszy, bo czas jest wielomianowy także względem $1/\eps$. Schematy aproksymacyjne są sposobem radzenia sobie z NP-trudnością, gdy optimum dokładne jest za drogie.

\subsection{Programowanie dynamiczne}

Programowanie dynamiczne jest techniką projektowania algorytmów dla problemów z optymalną podstrukturą i nakładającymi się podproblemami. Zamiast wielokrotnie rozwiązywać te same podproblemy, zapamiętuje się ich wyniki.

\subsubsection{Dwie podstawowe własności}

\begin{description}
  \item[Optymalna podstruktura] rozwiązanie optymalne problemu można zbudować z rozwiązań optymalnych podproblemów.
  \item[Nakładające się podproblemy] rekurencyjne rozwiązanie generowałoby te same podproblemy wielokrotnie.
\end{description}

Jeśli podproblemy się nie nakładają, zwykle wystarcza dziel i zwyciężaj, jak w mergesort. Jeśli nie ma optymalnej podstruktury, DP może dawać błędne rozwiązania.

\subsubsection{Top-down i bottom-up}

Podejście top-down z memoizacją zapisuje wyniki funkcji rekurencyjnej. Jest łatwe do napisania i oblicza tylko potrzebne stany.

Podejście bottom-up wypełnia tablicę w porządku zgodnym z zależnościami. Zwykle daje mniejszy narzut i łatwiej kontrolować pamięć.

\subsubsection{Schemat projektowania DP}

\begin{enumerate}
  \item Zdefiniować stan, czyli co opisuje podproblem.
  \item Określić wartość przechowywaną w stanie.
  \item Wyprowadzić rekurencję przejścia.
  \item Ustalić przypadki bazowe.
  \item Ustalić kolejność obliczeń.
  \item Odtworzyć rozwiązanie, jeśli potrzebna jest struktura, a nie tylko wartość.
  \item Oszacować liczbę stanów i koszt przejścia.
\end{enumerate}

\subsubsection{Przykład: najdłuższy wspólny podciąg}

Dla napisów $X[1..m]$ i $Y[1..n]$ definiujemy $dp[i][j]$ jako długość LCS prefiksów $X[1..i]$ i $Y[1..j]$:
\[
dp[i][j]=
\begin{cases}
0, & i=0\text{ lub }j=0,\\
dp[i-1][j-1]+1, & X_i=Y_j,\\
\max(dp[i-1][j],dp[i][j-1]), & X_i\ne Y_j.
\end{cases}
\]
Złożoność czasu i pamięci wynosi $\Theta(mn)$, choć pamięć dla samej długości można zredukować do $\Theta(\min(m,n))$.

\subsubsection{Przykład: plecak 0/1}

Dla przedmiotów o wagach $w_i$ i wartościach $v_i$ oraz pojemności $W$:
\[
dp[i][c]=\text{maksymalna wartość z pierwszych }i\text{ przedmiotów przy pojemności }c.
\]
Rekurencja:
\[
dp[i][c]=
\begin{cases}
dp[i-1][c], & w_i>c,\\
\max(dp[i-1][c],dp[i-1][c-w_i]+v_i), & w_i\le c.
\end{cases}
\]
Czas $\Theta(nW)$ jest pseudowielomianowy, bo zależy od wartości liczbowej $W$, a nie od liczby bitów jej zapisu.

\subsubsection{Przykład: Floyd-Warshall}

Dla najkrótszych ścieżek między wszystkimi parami:
\[
d^{(k)}_{ij}=\min\left(d^{(k-1)}_{ij},d^{(k-1)}_{ik}+d^{(k-1)}_{kj}\right).
\]
Stan oznacza najkrótszą ścieżkę z $i$ do $j$ używającą jako pośrednich tylko wierzchołków z $\{1,\dots,k\}$. Czas $\Theta(n^3)$, pamięć $\Theta(n^2)$.

\subsubsection{Przykład: Held-Karp dla TSP}

Stan:
\[
dp[S][j]=\text{koszt najkrótszej ścieżki startującej w }1,\text{ odwiedzającej }S,\text{ kończącej w }j.
\]
Rekurencja:
\[
dp[S][j]=\min_{i\in S, i\ne j} dp[S\setminus\{j\}][i]+c_{ij}.
\]
Czas $\Theta(n^2 2^n)$ jest dużo lepszy niż brute force, ale nadal wykładniczy.

\subsubsection{Pułapki}

Najczęstsze błędy:
\begin{itemize}
  \item stan nie zawiera całej informacji potrzebnej do przyszłych decyzji,
  \item rekurencja używa stanu, który nie został jeszcze obliczony,
  \item pomylenie indeksów i przypadków bazowych,
  \item nieświadome użycie pseudowielomianowej złożoności,
  \item odtworzenie rozwiązania nie jest zaplanowane,
  \item pamięć jest zbyt duża mimo akceptowalnego czasu.
\end{itemize}

\subsubsection{Co powiedzieć na obronie}

Programowanie dynamiczne stosuje się, gdy problem ma optymalną podstrukturę i nakładające się podproblemy. Definiujemy stan, rekurencję, przypadki bazowe i kolejność obliczeń. DP może być top-down z memoizacją albo bottom-up. Typowe przykłady to LCS, plecak, Floyd-Warshall, edycja napisów i Held-Karp dla TSP.
